\chapter{Aires, Volumes \& Conversions d'Unités}
\label{chap:aires_volumes}

% ==============================================================================
% 1. OBJECTIFS & COMPÉTENCES DU RÉFÉRENTIEL
% ==============================================================================
\begin{cadreretenu}[Objectifs \& Compétences du DNB Pro]
\begin{itemize}
    \item \textbf{Connaissances} : Formules d'aires usuelles (rectangle, triangle, disque, trapèze), formules de volumes (pavé droit, prisme droit, cylindre, cône, pyramide, boule), équivalence $1\text{ L} = 1\text{ dm}^3$.
    \item \textbf{Capacités} :
    \begin{itemize}
        \item \capRealiser\ Effectuer des conversions d'unités de longueurs, d'aires et de volumes/capacités.
        \item \capRealiser\ Calculer l'aire d'une surface complexe en la décomposant en figures simples.
        \item \capRealiser\ Calculer le volume d'un réservoir, d'un silo ou d'une cuve industrielle.
        \item \capAnalyser\ Déterminer une masse à partir d'un volume et d'une masse volumique ($\rho = \frac{m}{V}$).
    \end{itemize}
\end{itemize}
\end{cadreretenu}

\vspace{0.3cm}

% ==============================================================================
% 2. COURS : FORMULAIRE OFFICIEL AIRES ET VOLUMES
% ==============================================================================
\section{Cours : Formulaire d'Aires et de Volumes}

\subsection{Surfaces Planes Usuelles}

\noindent
\begin{tabularx}{\linewidth}{|X|X|X|}
\hline
\rowcolor{colDefBg} \textbf{Rectangle} & \textbf{Triangle} & \textbf{Disque} \\
\hline
\begin{center}
\begin{tikzpicture}[scale=0.6]
    \draw[fill=colDef!20, draw=colDef, thick] (0,0) rectangle (3,1.8);
    \node at (1.5,0.9) {$\mathcal{A} = L \times \ell$};
\end{tikzpicture}
\end{center} &
\begin{center}
\begin{tikzpicture}[scale=0.6]
    \draw[fill=colProp!20, draw=colProp, thick] (0,0) -- (3,0) -- (1,2) -- cycle;
    \draw[dashed] (1,2) -- (1,0);
    \node at (1.5,0.7) {$\mathcal{A} = \frac{b \times h}{2}$};
\end{tikzpicture}
\end{center} &
\begin{center}
\begin{tikzpicture}[scale=0.6]
    \draw[fill=colExem!20, draw=colExem, thick] (0,0) circle (1.2);
    \draw[->] (0,0) -- (1.2,0) node[midway, above, font=\tiny] {$R$};
    \node at (0,-0.3) [font=\footnotesize] {$\mathcal{A} = \pi R^2$};
\end{tikzpicture}
\end{center} \\
\hline
\end{tabularx}

\subsection{Solides et Volumes Usuels}

\begin{cadreprop}[Volumes des solides de l'espace]
\begin{itemize}
    \item \textbf{Pavé droit (Parallélépipède)} : $\mathcal{V} = L \times \ell \times h$
    \item \textbf{Prisme droit \& Cylindre de révolution} : $\mathcal{V} = \mathcal{A}_{\text{base}} \times h = \pi R^2 h$
    \item \textbf{Pyramide \& Cône de révolution} : $\mathcal{V} = \dfrac{1}{3} \times \mathcal{A}_{\text{base}} \times h = \dfrac{1}{3} \pi R^2 h$
    \item \textbf{Boule / Sphère} : $\mathcal{V} = \dfrac{4}{3} \pi R^3$ \quad et \quad $\mathcal{A}_{\text{sphère}} = 4 \pi R^2$
\end{itemize}
\end{cadreprop}

\begin{cadreretenu}[Conversions d'Unités Incontournables]
\[ 1\text{ m}^3 = 1\,000\text{ dm}^3 = 1\,000\text{ L} \qquad\text{et}\qquad 1\text{ dm}^3 = 1\text{ L} = 1\,000\text{ cm}^3 = 1\,000\text{ mL} \]
\end{cadreretenu}

% ==============================================================================
% 3. QUESTIONS FLASH / AUTOMATISMES
% ==============================================================================
\section{Questions Flash / Automatismes}

\begin{automatisme}[Conversions \& Formules (10 min)]
\noindent
\begin{tabularx}{\linewidth}{X|X}
\textbf{1.} $2{,}5\text{ m}^3 = \pointille[1.5cm]\text{ L}$ & \textbf{4.} Volume d'un cube d'arête $a = 4\text{ cm}$ : $V = \pointille[1.5cm]\text{ cm}^3$ \\
\textbf{2.} $750\text{ cm}^3 = \pointille[1.5cm]\text{ L}$ & \textbf{5.} Aire d'un disque de rayon $R = 10\text{ cm}$ : $\mathcal{A} \approx \pointille[1.5cm]\text{ cm}^2$ \\
\textbf{3.} $12\text{ m}^2 = \pointille[1.5cm]\text{ dm}^2$ & \textbf{6.} Masse volumique de l'eau : $1\text{ L d'eau} = \pointille[1cm]\text{ kg}$ \\
\end{tabularx}
\end{automatisme}

% ==============================================================================
% 4. TRAVAUX DIRIGÉS (TD)
% ==============================================================================
\section{Travaux Dirigés (TD) : Exercices d'Entraînement}

\begin{exercice}[4 pts]{\capRealiser}{Cuve de Récupération d'Eau de Pluie}
Une cuve cylindrique de récupération d'eau a un diamètre intérieur de $D = 1{,}20\text{ m}$ et une hauteur utile de $h = 1{,}80\text{ m}$.
\begin{enumerate}
    \item Calculer le rayon intérieur $R$ en mètres.
    \item Calculer le volume de la cuve en $\text{m}^3$. Arrondir au centième.
    \item En déduire la capacité maximale de la cuve en Litres.
\end{enumerate}
\lignesreponse[4]
\end{exercice}

\begin{exercice}[4 pts]{\capAnalyser}{Coulage d'une Dalle en Béton (BTP)}
Un maçon réalise une dalle en béton rectangulaire de longueur $L = 8\text{ m}$, de largeur $\ell = 4{,}50\text{ m}$ et d'épaisseur $e = 15\text{ cm}$.
\begin{enumerate}
    \item Exprimer l'épaisseur $e$ en mètres.
    \item Calculer le volume de béton nécessaire en $\text{m}^3$.
    \item Sachant que la masse volumique du béton armé est $\rho = 2\,400\text{ kg/m}^3$, calculer la masse totale de la dalle en tonnes.
\end{enumerate}
\lignesreponse[5]
\end{exercice}

% ==============================================================================
% 5. VERS LE BREVET (DNB PRO)
% ==============================================================================
\section{Vers le Brevet (DNB Pro)}

\begin{sujetdnb}[DNB Pro Métropole][9 pts]{Traitement d'une Piscine Municipale}
Une piscine municipale dispose d'un bassin de forme pavé droit de longueur $25\text{ m}$, de largeur $10\text{ m}$ et de profondeur constante $1{,}60\text{ m}$.

\begin{enumerate}
    \item \capRealiser\ Calculer le volume d'eau contenu dans ce bassin en $\text{m}^3$, puis en Litres.
    \item \capAnalyser\ Pour traiter l'eau, le maître-nageur doit verser $2\text{ cL}$ de produit désinfectant pour $1\,000\text{ L}$ d'eau.
    Calculer la quantité totale de désinfectant à verser (en Litres).
    \item \capValider\ La pompe de filtration a un débit de $40\text{ m}^3/\text{h}$. Combien d'heures faut-il pour renouveler la totalité de l'eau du bassin ?
\end{enumerate}
\lignesreponse[6]
\end{sujetdnb}

% ==============================================================================
% 6. CORRIGÉS DÉTAILLÉS
% ==============================================================================
\section{Corrigés Détaillés du Chapitre}

\begin{solution}[Corrigé des Exercices et du Sujet DNB]
\textbf{Exercice 1 (Cuve)} :
1. $R = 1{,}20 / 2 = \repbox{0{,}60\text{ m}}$. \\
2. $V = \pi \times R^2 \times h = \pi \times 0{,}60^2 \times 1{,}80 \approx \repbox{2{,}04\text{ m}^3}$. \\
3. $1\text{ m}^3 = 1\,000\text{ L} \implies 2{,}04 \times 1\,000 = \repbox{2\,040\text{ Litres}}$.

\vspace{2mm}
\textbf{Exercice 2 (Béton BTP)} :
1. $e = 15\text{ cm} = \repbox{0{,}15\text{ m}}$. \\
2. $V = 8 \times 4{,}50 \times 0{,}15 = \repbox{5{,}40\text{ m}^3}$. \\
3. Masse $= 5{,}40 \times 2\,400 = 12\,960\text{ kg} = \repbox{12{,}96\text{ tonnes}}$.

\vspace{2mm}
\textbf{Sujet DNB Pro (Piscine)} :
1. $V = 25 \times 10 \times 1{,}60 = \repbox{400\text{ m}^3} = \repbox{400\,000\text{ Litres}}$. \\
2. Produit $= \frac{400\,000}{1\,000} \times 2\text{ cL} = 400 \times 2 = 800\text{ cL} = \repbox{8\text{ Litres}}$. \\
3. Temps de filtration $= \frac{400}{40} = \repbox{10\text{ heures}}$.
\end{solution}
